\documentclass[11pt]{article}
\usepackage[margin=1in]{geometry}
\usepackage{amsmath,amssymb,amsthm}
\usepackage{graphicx}
\usepackage[hidelinks]{hyperref}

\newtheorem{theorem}{Theorem}
\newtheorem{lemma}{Lemma}
\newtheorem{corollary}{Corollary}
\newtheorem{remark}{Remark}

\newcommand{\N}{\mathbb N}
\newcommand{\R}{\mathbb R}
\newcommand{\eps}{\varepsilon}
\newcommand{\norm}[1]{\left\lVert #1 \right\rVert}
\newcommand{\bbeta}{\bar{\beta}}

\title{A Counterexample to Extending One-Point Asymptotic Beta Positivity Beyond One}
\author{Run \texttt{fa\_banach\_001}, lane 3}
\date{June 26, 2026}

\begin{document}
\maketitle

\section*{Source Question}

The source paper is F. P. Baudier and G. Lancien,
\emph{An asymptotic analog of a local-to-global phenomenon for uniformly convex
renormings}, arXiv:2401.17465.  Theorem 1.2 states that if a Banach space
$(X,\norm{\cdot})$ satisfies
\[
  \bbeta_{\norm{\cdot}}(t_0)>0
  \quad\text{for some }t_0\in(0,1),
\]
then $X$ admits an equivalent norm with Rolewicz's property $(\beta)$.
In the final paragraph, the authors observe that the asymptotic beta modulus
also makes sense on the larger interval $[0,K(X))$, where
\[
  K(X)=\sup\{\lambda>0:\text{ there is a sequence }(x_n)\subset B_X
       \text{ with }\inf_{n\ne m}\norm{x_n-x_m}\ge \lambda\},
\]
and ask:

\begin{quote}
It would be interesting to see if the condition in Theorem 1.2 could be pushed
beyond 1.
\end{quote}

\begin{center}
\includegraphics[width=\textwidth]{figures/open_problem_crop.png}
\end{center}

\section*{Answer}

The answer to the natural strengthened statement is negative.  There is a
non-reflexive Banach space whose given norm has positive asymptotic beta
modulus at every scale strictly larger than $1$ up to its Kottman range.
Consequently, one-point positivity at some $t_0\in(1,K(X))$ does not force the
space to admit an equivalent property-$(\beta)$ norm.

\section*{Definitions}

For a Banach space $(X,\norm{\cdot})$ and a parameter for which separated
sequences exist in the unit ball, use the extended asymptotic beta modulus
\[
\bbeta_X(t)=
\inf \left\{
 \sup_{n\in\N}\left(1-\norm{\frac{x+y_n}{2}}\right):
 x\in B_X,\ (y_n)\subset B_X,\
 \inf_{n\ne m}\norm{y_n-y_m}\ge t
\right\}.
\]
Rolewicz's property $(\beta)$ is equivalent to positivity of this modulus on
$(0,1]$ for the given norm.  It is standard, and also recalled in the source
paper, that property $(\beta)$ implies reflexivity.  Hence a non-reflexive
space cannot admit any equivalent norm with property $(\beta)$.

\section*{Proof Intuition}

Let
\[
  X=\ell_2\oplus_2 c_0,
  \qquad
  \norm{(h,z)}=\big(\norm{h}_2^2+\norm{z}_\infty^2\big)^{1/2}.
\]
The $c_0$ summand makes $X$ non-reflexive.  However, $c_0$ by itself cannot
support infinite separated sequences at any scale above $1$: in any bounded
sequence in $c_0$, two terms are within $1+\eps$ in sup norm.  Thus any
sequence in $X$ separated by a scale $t>1$ must, for a suitable pair of terms,
use a definite amount of separation in the Hilbert coordinate.  The Hilbert
coordinate is strictly convex enough that two such Hilbert-separated terms
cannot both have almost norm-one midpoint with the same point $x\in B_X$.

\section*{Auxiliary Lemmas}

\begin{lemma}[The $c_0$ Kottman threshold]\label{lem:c0}
Let $(z_n)$ be a sequence in $B_{c_0}$.  For every $c>1$ there exist distinct
indices $n,m$ such that
\[
  \norm{z_n-z_m}_\infty\le c .
\]
\end{lemma}

\begin{proof}
Write $c=1+\eps$ with $\eps>0$.  Since $c_0\subset[-1,1]^\N$, a diagonal
argument gives a subsequence, still denoted $(z_n)$, which converges
coordinatewise.  Choose a finite set $F\subset\N$ such that
$|z_1(k)|\le \eps/2$ for $k\notin F$.  By coordinatewise convergence, pick
$m>1$ so large that $|z_m(k)-z_1(k)|\le \eps$ for all $k\in F$.  If
$k\notin F$, then
\[
  |z_m(k)-z_1(k)|\le |z_m(k)|+|z_1(k)|\le 1+\eps/2<1+\eps.
\]
Thus $\norm{z_m-z_1}_\infty\le 1+\eps=c$.
\end{proof}

\begin{lemma}[Two-point Hilbert-sum midpoint drop]\label{lem:twopoint}
Let $H$ be a Hilbert space, let $Z$ be a Banach space, and put
$E=H\oplus_2 Z$.  If $0\le c<t\le 2$, then there is a number
$\eta=\eta(c,t)>0$ such that the following holds.  Whenever
\[
  x=(u,v),\quad y=(h,z),\quad y'=(h',z')\in B_E
\]
satisfy
\[
  \norm{z-z'}_Z\le c
  \quad\text{and}\quad
  \norm{y-y'}_E\ge t,
\]
then
\[
  \max\left\{
  1-\norm{\frac{x+y}{2}}_E,\,
  1-\norm{\frac{x+y'}{2}}_E
  \right\}\ge \eta .
\]
\end{lemma}

\begin{proof}
Assume the contrary.  Then there are triples
$x_j=(u_j,v_j)$, $y_j=(h_j,z_j)$ and
$y'_j=(h'_j,z'_j)$ in $B_E$ such that
\[
 \norm{z_j-z'_j}_Z\le c,\qquad
 \norm{y_j-y'_j}_E\ge t,
\]
while
\[
  \norm{\frac{x_j+y_j}{2}}_E\to 1
  \quad\text{and}\quad
  \norm{\frac{x_j+y'_j}{2}}_E\to 1 .
\]
Equivalently, $\norm{x_j+y_j}_E\to2$ and
$\norm{x_j+y'_j}_E\to2$.

From $\norm{x_j+y_j}_E\le\norm{x_j}_E+\norm{y_j}_E\le2$, we get
$\norm{x_j}_E\to1$ and $\norm{y_j}_E\to1$; similarly
$\norm{y'_j}_E\to1$.  Put
\[
 A_j=(\norm{u_j}_H,\norm{v_j}_Z),\quad
 B_j=(\norm{h_j}_H,\norm{z_j}_Z),\quad
 C_j=(\norm{h'_j}_H,\norm{z'_j}_Z)
\]
in the Euclidean plane.  Since
\[
\norm{x_j+y_j}_E
\le \norm{A_j+B_j}_2
\le \norm{A_j}_2+\norm{B_j}_2
\le 2
\]
and the left side tends to $2$, equality is asymptotically attained in the
Euclidean triangle inequality.  Hence $\norm{A_j-B_j}_2\to0$.  The same
argument with $y'_j$ gives $\norm{A_j-C_j}_2\to0$.

Moreover,
\[
0\le
(\norm{u_j}_H+\norm{h_j}_H)^2-\norm{u_j+h_j}_H^2
\le
\norm{A_j+B_j}_2^2-\norm{x_j+y_j}_E^2\to0.
\]
Using the Hilbert identity
\[
\norm{u_j-h_j}_H^2
 =
(\norm{u_j}_H-\norm{h_j}_H)^2
+\big((\norm{u_j}_H+\norm{h_j}_H)^2-\norm{u_j+h_j}_H^2\big),
\]
we obtain $\norm{u_j-h_j}_H\to0$.  Repeating the same argument for $y'_j$
gives $\norm{u_j-h'_j}_H\to0$, and therefore
$\norm{h_j-h'_j}_H\to0$.

It follows that
\[
\limsup_j \norm{y_j-y'_j}_E^2
 =
\limsup_j\left(\norm{h_j-h'_j}_H^2+\norm{z_j-z'_j}_Z^2\right)
\le c^2<t^2,
\]
contradicting $\norm{y_j-y'_j}_E\ge t$ for all $j$.
\end{proof}

\section*{Main Counterexample}

\begin{theorem}\label{thm:counterexample}
Let $X=\ell_2\oplus_2 c_0$ with its natural $\ell_2$-sum norm.  Then:
\begin{enumerate}
\item $K(X)\ge\sqrt2$.
\item $\bbeta_X(t)>0$ for every $t\in(1,\sqrt2)$.
\item $X$ does not admit an equivalent norm with property $(\beta)$.
\end{enumerate}
In particular, $X$ satisfies $\bbeta_X(t_0)>0$ for some
$t_0\in(1,K(X))$ but is not property-$(\beta)$ renormable.
\end{theorem}

\begin{proof}
For (1), let $(e_n)$ be an orthonormal sequence in $\ell_2$.  The points
$(e_n,0)$ lie in $B_X$ and satisfy
\[
  \norm{(e_n,0)-(e_m,0)}_X=\sqrt2
  \qquad(n\ne m).
\]
Thus $K(X)\ge\sqrt2$.

For (2), fix $t\in(1,\sqrt2)$ and choose $c$ with $1<c<t$, for instance
$c=(1+t)/2$.  Let $\eta=\eta(c,t)>0$ be given by Lemma~\ref{lem:twopoint}.
We prove that $\bbeta_X(t)\ge\eta$.

Take an arbitrary $x=(u,v)\in B_X$ and an arbitrary sequence
$y_n=(h_n,z_n)\in B_X$ with
\[
  \inf_{n\ne m}\norm{y_n-y_m}_X\ge t.
\]
By Lemma~\ref{lem:c0}, there are distinct indices $n,m$ such that
\[
  \norm{z_n-z_m}_\infty\le c.
\]
Since also $\norm{y_n-y_m}_X\ge t$, Lemma~\ref{lem:twopoint} applied in
$H=\ell_2$ and $Z=c_0$ yields
\[
  \max\left\{
  1-\norm{\frac{x+y_n}{2}}_X,\,
  1-\norm{\frac{x+y_m}{2}}_X
  \right\}\ge \eta .
\]
Consequently
\[
  \sup_k\left(1-\norm{\frac{x+y_k}{2}}_X\right)\ge \eta .
\]
Because $x$ and the separated sequence $(y_n)$ were arbitrary, the defining
infimum gives $\bbeta_X(t)\ge\eta>0$.

For (3), $X$ contains the closed subspace $\{0\}\oplus c_0$, so $X$ is
non-reflexive.  If $X$ admitted an equivalent norm with property $(\beta)$,
then $X$ would be reflexive, since property $(\beta)$ implies reflexivity.
This contradiction proves that no equivalent property-$(\beta)$ norm exists.
\end{proof}

\section*{Verification Notes}

The proof is deterministic and uses no numerical computation.  The two points
where human review should focus are:
\begin{itemize}
\item Lemma~\ref{lem:c0}, which records exactly that $c_0$ has no infinite
unit-ball separated sequence at a scale strictly larger than $1$.
\item Lemma~\ref{lem:twopoint}, where near norm-one midpoints in
$H\oplus_2 Z$ force the Hilbert parts of both points to converge to the
Hilbert part of the common base point.
\end{itemize}
Together these lemmas give a uniform midpoint drop for every $t>1$ and hence
positive asymptotic beta modulus at any chosen $t_0\in(1,\sqrt2)$.

\section*{Novelty and Search Bounds}

Before promotion, the run indexes were searched for arXiv:2401.17465, for the
phrases \texttt{Kottman}, \texttt{pushed beyond 1}, and
\texttt{property beta}, and for the earlier local attempt on the
Kottman-domain beta extension.  No durable solution packet or proof-gap packet
for this target was found; the previous failed attempt is copied into this
packet under \texttt{history/}.

Bounded web searches on June 26, 2026 for
\texttt{"2401.17465" "pushed beyond 1"},
\texttt{"property beta" "Kottman's constant"},
\texttt{"Rolewicz" "Kottman" "beta" "c0"},
and the exact paper title returned the source arXiv page and no explicit later
answer to the Kottman-domain question.  Novelty confidence is therefore
moderate: the construction is short and should be checked against the
specialist renorming literature, but no known answer was found in the bounded
search.

\section*{Human Review Recommendation}

Treat this as a likely valid counterexample to the strengthened
Kottman-domain version of Theorem 1.2.  A reviewer should verify that the
source question is intended to include the natural formulation ``some
$t_0\in(1,K(X))$'' and should audit Lemma~\ref{lem:twopoint}'s compactness
argument.  If both checks pass, the example $X=\ell_2\oplus_2 c_0$ gives a
negative answer.

\begin{thebibliography}{9}

\bibitem{BaudierLancien}
F. P. Baudier and G. Lancien,
\emph{An asymptotic analog of a local-to-global phenomenon for uniformly
convex renormings}, arXiv:2401.17465.

\bibitem{Rolewicz}
S. Rolewicz,
\emph{On drop property}, Studia Math. 85 (1987), 27--35.

\bibitem{Kutzarova}
D. Kutzarova,
\emph{$k$-$\beta$ and $k$-nearly uniformly convex Banach spaces},
J. Math. Anal. Appl. 162 (1991), no. 2, 322--338.

\bibitem{EltonOdell}
J. Elton and E. Odell,
\emph{The unit ball of every infinite-dimensional normed linear space contains
a $(1+\epsilon)$-separated sequence}, Colloq. Math. 44 (1981), 105--109.

\end{thebibliography}

\end{document}
